This paper presented a dynamic OSM-based route-candidate reranking system that builds pseudo-history profiles from public GPS-derived movement signals, extracts interpretable OSM route features, and reranks generated route candidates using profile, prompt, or hybrid modes. The main methodological point is separation: reconstruction quality, candidate-generation quality, and reranking quality are evaluated independently so that noisy public traces are not mistaken for clean user histories.

The current results show a clear research path without claiming final superiority over learned route-search systems. In the Toronto OSM public-trace setting, the pipeline produces usable exploratory pseudo-history signals and reports reconstruction diagnostics that expose route-fidelity gaps. In the primary Porto public-dataset benchmark, a trajectory-derived vehicle profile improves over shortest-distance and a generic profile on ranking metrics, while path metrics identify map matching and candidate generation as the next technical leverage points. The result is a reproducible research framework with a concrete public-benchmark agenda: larger-scale Porto evaluation, reproduced external baselines such as NASR under matched preprocessing and splits, and stronger map matching can now be added without changing the system architecture or evaluation philosophy.
